\chapter{Lattice Reflection Positivity}
\label{ch:lattice-reflection-positivity}

\ConventionsEuclidean

Reflection positivity is the lattice analogue of the positivity axiom used in
Euclidean reconstruction.  It is a property of a regulated measure and a
chosen reflection, not an automatic property of every discretization.  This
chapter states the condition and verifies it for the basic nearest-neighbor
scalar lattice action.

\section{Reflection and Positive Algebra}

Let \(\Lambda\subset a\mathbb Z^D\) be a finite reflection-invariant lattice
with reflection \(\vartheta\) across the hyperplane \(x_0=0\).  Let
\(\phi_x\) be a real scalar variable at \(x\in\Lambda\).  Let
\(\mathcal A_+\) be the algebra of polynomial functions of the variables with
\(x_0>0\).  The reflection acts antilinearly on functions by
\[
  (\Theta F)(\phi)=\overline{F(\phi_{\vartheta x})}.
\]
For a probability measure \(d\mu(\phi)\), define
\[
  \langle F\rangle=\int F(\phi)\,d\mu(\phi).
\]

\begin{definition}[Lattice reflection positivity]
\label{def:lattice-reflection-positivity}
The measure \(d\mu\) is reflection positive with respect to
\(\vartheta\) if
\[
  \langle \Theta F\,F\rangle\ge 0
\]
for every \(F\in\mathcal A_+\).
\end{definition}

\section{Nearest-Neighbor Scalar Measure}

Consider the finite-volume lattice measure
\[
  d\mu(\phi)=Z^{-1}e^{-S(\phi)}\prod_{x\in\Lambda}d\phi_x
\]
with
\[
  S(\phi)=
  \frac12\sum_{\langle x y\rangle}c_{xy}(\phi_x-\phi_y)^2
  +\sum_x V(\phi_x),
  \qquad c_{xy}>0,
\]
where \(V\) is a real potential for which the finite-volume integral exists.
Split the action as
\[
  S=S_+ + \Theta S_+ + S_0
\]
where \(S_+\) contains terms strictly in the positive half, \(\Theta S_+\)
its reflected copy, and \(S_0\) the terms on or crossing the reflection
plane.  The crossing terms have the form
\[
  \frac12 c(\phi_x-\phi_{\vartheta x})^2
  =
  \frac12c\phi_x^2+\frac12c\phi_{\vartheta x}^2
  -c\phi_x\phi_{\vartheta x}.
\]
The first two terms can be assigned to \(S_+\) and \(\Theta S_+\).  The
coupling across the plane contributes the positive kernel
\[
  \exp(c\phi_x\phi_{\vartheta x})
  =
  \sum_{n=0}^{\infty}\frac{c^n}{n!}\phi_x^n\phi_{\vartheta x}^n .
\]

\begin{proposition}[Scalar lattice reflection positivity]
\label{prop:scalar-lattice-reflection-positivity}
The nearest-neighbor scalar lattice measure above is reflection positive.
\end{proposition}

\begin{proof}
After assigning all non-crossing and square terms to the two halves, the only
factor coupling reflected variables is a product of kernels
\(\exp(c_i\phi_i\phi_{\vartheta i})\) with \(c_i>0\).  Expanding each such
kernel in its power series gives a sum of terms
\[
  a_{\mathbf n}\,
  M_{\mathbf n}(\phi_+)\,
  \overline{M_{\mathbf n}(\phi_+)}
  \]
after reflection, with coefficients \(a_{\mathbf n}\ge 0\).  For
\(F\in\mathcal A_+\), the integral \(\langle \Theta F\,F\rangle\) becomes a
sum of absolute squares of half-lattice integrals weighted by
\(a_{\mathbf n}\).  Every term is nonnegative, hence the sum is nonnegative.
\end{proof}

\section{Transfer Matrix Interpretation}

When the lattice is periodic in Euclidean time, reflection positivity permits
the construction of a pre-Hilbert space from \(\mathcal A_+\) by quotienting
the null space of the sesquilinear form
\[
  (F,G)=\langle \Theta F\,G\rangle .
\]
Time translation by one lattice spacing acts as a contraction on this
pre-Hilbert space.  Under the additional hypotheses required for a transfer
matrix, this contraction is \(e^{-aH}\) for a positive self-adjoint lattice
Hamiltonian \(H\).  This is the finite-regulator version of the positivity
step in OS reconstruction.

\section{Gauge Fields}

For lattice gauge theory the variables are group elements \(U_\ell\) on
oriented links.  Reflection positivity of Wilson's plaquette action depends
on assigning plaquettes to the two halves and expanding the cross-plane
plaquette weights in positive characters.  The property is sensitive to the
choice of lattice action; improved actions or gauge-fixing terms require a
separate positivity check.

\begin{openproblem}[Reflection-positive continuum limits]
\label{op:reflection-positive-continuum-limits}
Give verifiable criteria ensuring that a sequence of reflection-positive
lattice measures converges to OS-positive continuum Schwinger functions with
nontrivial local fields.  The criteria must include tightness or correlation
bounds, control of counterterms, and compatibility of reflection with the
continuum scaling limit.
\end{openproblem}
